% Directly inputtable paper section. Author: ChatGPT.
\section{Quadratic Veronese peeling and valuation-layer flags}
\label{sec:quadratic-veronese-peeling}

\noindent\textbf{Author: ChatGPT.}
This section isolates the specialization $y=x^2$ of the five-term contact
relation, transports it to primitive integer triples, and audits the
resulting one-move policy on a full-premise Pythagorean-square family.  The
fixed policy fails three necessary uniform estimates.  The unrestricted
multi-move contact-complex route remains open.

For $z\in\mathbf Q\setminus\{0,1\}$, let
\[
 \Omega(z)=d(z)\wedge d(1-z),
 \qquad d:\mathbf Q^\times\longrightarrow
   \bigoplus_p\mathbf Z[p].
\]
Recall the five-term identity
\begin{equation}
 \Omega(x)-\Omega(y)+\Omega(y/x)
 -\Omega\!\left(\frac{y(1-x)}{x(1-y)}\right)
 +\Omega\!\left(\frac{1-x}{1-y}\right)=0.
 \label{eq:qvp-five-term}
\end{equation}

\begin{theorem}[Quadratic Veronese peeling]
\label{thm:qvp-peeling}
The specialization $y=x^2$ is admissible in
\eqref{eq:qvp-five-term} exactly when
$x\notin\{0,1,-1\}$.  On this domain,
\begin{equation}
 \Omega(x^2)=2\Omega(x)
 -2\Omega\!\left(\frac{x}{1+x}\right).
 \label{eq:qvp-peeling}
\end{equation}
\end{theorem}

\begin{proof}
The five-term domain says that $x,x^2$ are nonzero, different from one, and
different from each other.  Since
$x^2-1=(x-1)(x+1)$ and $x^2-x=x(x-1)$, this is equivalent to the three
stated exclusions.  Direct field algebra then gives
\[
 \frac{x^2}{x}=x,
 \quad
 \frac{x^2(1-x)}{x(1-x^2)}=\frac{x}{1+x},
 \quad
 \frac{1-x}{1-x^2}=\frac1{1+x}.
\]
The final two arguments are complementary.  Alternation gives
$\Omega(1-z)=-\Omega(z)$, so substitution in
\eqref{eq:qvp-five-term} proves \eqref{eq:qvp-peeling}.  No sign condition on
$x$ is required, since $d(-1)=0$ for finite-prime divisors.
\end{proof}

Writing $X=d(x)$, $U=d(1-x)$, and $W=d(1+x)$ exposes the coefficientwise
identity behind the proof:
\begin{equation}
 (2X)\wedge(U+W)
 =2X\wedge U-2(X-W)\wedge(-W).
 \label{eq:qvp-vector}
\end{equation}

\subsection{Primitive integer transport}

Let $(a,b,c)$ be a positive pairwise-coprime triple with $a+b=c$ and put
$x=a/c$.

\begin{proposition}[Quadratic transform]
\label{prop:qvp-integer-transform}
The triple
\[
 T_2(a,b,c)=\bigl(a^2,b(a+c),c^2\bigr)
\]
is positive, pairwise coprime, and satisfies
$a^2+b(a+c)=c^2$.  Moreover,
\begin{equation}
 \Omega\bigl(a^2,b(a+c),c^2\bigr)
 =2\Omega(a,b,c)-2\Omega(a,c,a+c).
 \label{eq:qvp-integer}
\end{equation}
\end{proposition}

\begin{proof}
The sum identity follows from $b=c-a$:
$a^2+(c-a)(a+c)=c^2$.  A common prime of $a^2$ and $b(a+c)$ would divide
$a$ and either $b$ or $a+c$, contradicting the original pairwise
coprimality.  The same argument applies to $b(a+c)$ and $c^2$, while
$a^2$ and $c^2$ are coprime directly.  Equation
\eqref{eq:qvp-integer} is Theorem~\ref{thm:qvp-peeling} with $x=a/c$.
\end{proof}

For $t\geq1$, set
\[
 X=2t+1,\qquad Y=2t(t+1),\qquad Z=2t^2+2t+1.
\]
Then
\[
 X^2+Y^2=Z^2,\qquad Z-Y=1,\qquad Y+Z=X^2,
\]
and $X,Y,Z$ are pairwise coprime.  Taking $x=Y/Z$ in
\eqref{eq:qvp-peeling} yields the full-premise identity
\begin{equation}
 \boxed{\displaystyle
 \Omega(Y^2,X^2,Z^2)
 =2\Omega(Y,1,Z)-2\Omega(Y,Z,X^2).\ }
 \label{eq:qvp-pythagorean}
\end{equation}
Thus the first auxiliary cell has no nontrivial middle divisor, and only the
third leg of the second auxiliary cell retains a visible square.

\subsection{Exact positive-cost redistribution}

Put
\[
 u=\log X,\quad v=\log Y,\quad w=\log Z,
\]
and
\[
 \beta=\log\operatorname{rad}(X),\quad
 \alpha=\log\operatorname{rad}(Y),\quad
 \gamma=\log\operatorname{rad}(Z).
\]
Write
$\delta_X=u-\beta$, $\delta_Y=v-\alpha$, and
$\delta_Z=w-\gamma$.  For the declared outer-square decomposition, the
full, coherent, and residual costs of the squared cell are
\begin{align}
 \Phi_{\mathrm{sq}}&=4(uv+uw+vw),\nonumber\\
 V_{\mathrm{sq}}&=\Phi_{\mathrm{sq}},\nonumber\\
 R_{\mathrm{sq}}&=2\{\delta_Y(u+w)+\delta_X(v+w)
                         +\delta_Z(u+v)\}.
 \label{eq:qvp-square-costs}
\end{align}
After weighting both right-hand cells in
\eqref{eq:qvp-pythagorean} by the absolute coefficient two, their costs are
\begin{align}
 \Phi_{\mathrm{peel}}&=2vw+2\{vw+2u(v+w)\},\nonumber\\
 V_{\mathrm{peel}}&=2u(v+w),\nonumber\\
 R_{\mathrm{peel}}&=4\delta_Y(u+w)+2\delta_X(v+w)
                         +4\delta_Z(u+v).
 \label{eq:qvp-peeled-costs}
\end{align}

\begin{proposition}[Conservation and redistribution]
\label{prop:qvp-cost-redistribution}
One has
\begin{align}
 \Phi_{\mathrm{peel}}&=\Phi_{\mathrm{sq}},
 \label{eq:qvp-area-conservation}\\
 V_{\mathrm{peel}}&<\tfrac12V_{\mathrm{sq}},
 \label{eq:qvp-veronese-drop}\\
 R_{\mathrm{peel}}-R_{\mathrm{sq}}
 &=2\delta_Y(u+w)+2\delta_Z(u+v)\geq0.
 \label{eq:qvp-residual-rise}
\end{align}
In addition,
$V_{\mathrm{peel}}/V_{\mathrm{sq}}\to1/4$ as $t\to\infty$.
\end{proposition}

\begin{proof}
Expansion of \eqref{eq:qvp-peeled-costs} proves
\eqref{eq:qvp-area-conservation}.  The difference between half the old
coherent cost and the new coherent cost is $2vw>0$.  Direct subtraction
proves \eqref{eq:qvp-residual-rise}.  Finally
$u\sim\log t$ and $v,w\sim2\log t$, which gives the stated limit.
\end{proof}

For calibrated cost $Q=M+V=2\Phi-R$, the same computation gives
\begin{equation}
 Q_{\mathrm{peel}}-Q_{\mathrm{sq}}
 =-2\delta_Y(u+w)-2\delta_Z(u+v).
 \label{eq:qvp-q-drop}
\end{equation}
The coherent improvement has therefore been transferred to residual cost;
the total $Q+R=2\Phi$ has not decreased.

\subsection{The primewise valuation-layer flag}

For $k\geq1$, define
\[
 R_k(n)=\prod_{p:\,v_p(n)\geq k}p,
 \qquad \lambda_k(n)=\log R_k(n).
\]
Then
\begin{equation}
 \log n=\sum_{k\geq1}\lambda_k(n),\qquad
 \log\operatorname{rad}(n)=\lambda_1(n),\qquad
 \log\frac{n}{\operatorname{rad}(n)}
 =\sum_{k\geq2}\lambda_k(n).
 \label{eq:qvp-layer-cake}
\end{equation}
This filtration sees repeated prime mass even when the gcd of all exponents
is one; $2s^2$ with $s$ odd is the basic example.  Bilinearity gives the
exact layer-pair expansion
\begin{equation}
 \Phi(a,b,c)=\sum_{k,l\geq1}
 \bigl(\lambda_k(a)\lambda_l(b)
      +\lambda_k(b)\lambda_l(c)
      +\lambda_k(c)\lambda_l(a)\bigr).
 \label{eq:qvp-layer-pairs}
\end{equation}
Squaring duplicates every layer:
\begin{equation}
 R_{2j-1}(n^2)=R_{2j}(n^2)=R_j(n).
 \label{eq:qvp-layer-duplication}
\end{equation}
Consequently the square-cell layer-pair area is replicated four times, and
the layer ledger recovers the conservation law
\eqref{eq:qvp-area-conservation}.

The flag also rules out a naive uniform depth cutoff.  With $t=2^m$,
$v_2(Y^2)=2m+2$, so the prime $2$ occupies arbitrarily high layers inside
the genuine primitive cells $(Y^2,X^2,Z^2)$.  Weighted layer decay and
primewise matching across several five-term moves are not refuted by this
example.

\subsection{Three fixed-policy no-go theorems}

The first no-go is immediate from
\eqref{eq:qvp-area-conservation}.  If some $\eta>0$ and $K$ satisfied
\[
 \Phi_{\mathrm{peel}}
 \leq(1-\eta)\Phi_{\mathrm{sq}}+K(2w)
\]
for all $t$, then $\eta\Phi_{\mathrm{sq}}\leq2Kw$.  But $Z=Y+1\leq2Y$
and $Y\geq4$, so $w\leq3v/2$ and
\[
 \frac{\Phi_{\mathrm{sq}}}{2w}\geq\frac43u\longrightarrow\infty.
\]
Thus no positive one-step area contraction exists.

For the calibrated-boundary coefficient, expansion gives
\begin{align}
 Q_{\mathrm{peel}}={}&2u(v+w)+4\alpha(u+w)+4\gamma(u+v)
                         +2\beta(v+w),
 \label{eq:qvp-q-formula}\\
 Q_{\mathrm{peel}}-2H\rho={}&2u(v+w)+4\alpha u
   +4\gamma(u+v-w)+2\beta(v-w),
 \label{eq:qvp-boundary-excess}
\end{align}
where $H=2w$ and $\rho=\alpha+\beta+\gamma$.  Since $Y<Z\leq XY$,
$0\leq\beta\leq u$, and $\alpha,\gamma\geq0$, the last display implies
\begin{equation}
 Q_{\mathrm{peel}}-2H\rho\geq4uv.
 \label{eq:qvp-boundary-lower}
\end{equation}
Dividing by $H$ and using $w\leq3v/2$ shows that the excess is an unbounded
multiple of height.  Hence the fixed chain cannot satisfy
$Q_{\mathrm{peel}}\leq2H\rho+LH$ for any constant $L$.

Finally, the policy requiring the following bound for every
$\varepsilon>0$ is false:
\[
 R_{\mathrm{peel}}\leq
 \varepsilon Q_{\mathrm{peel}}+K_\varepsilon H
\]
with a uniform $K_\varepsilon$.  Set $t=2^m$ and
$L_m=\log t$.  For $m\geq2$,
\[
 L_m\leq u\leq2L_m,
 \quad 2L_m\leq v,w\leq3L_m,
 \quad \delta_Y\geq L_m.
\]
Thus
\[
 R_{\mathrm{peel}}\geq12L_m^2,
 \qquad \Phi_{\mathrm{sq}}\leq84L_m^2,
 \qquad H\leq6L_m.
\]
As $Q_{\mathrm{peel}}\leq2\Phi_{\mathrm{sq}}$, at
$\varepsilon=1/100$ one obtains
\[
 R_{\mathrm{peel}}-\frac1{100}Q_{\mathrm{peel}}
 \geq\frac{258}{25}L_m^2,
\]
which cannot be absorbed into any uniform $K_{1/100}H$.  Thus the single
admissible value $\varepsilon=1/100$ already refutes the universal-in-
$\varepsilon$ policy; no assertion is made here that every positive
$\varepsilon$ fails.  This theorem concerns the
declared outer-square split.  It does not claim that the same subsequence
refutes the maximal common-exponent residual.

All triples used above satisfy positivity, the sum equation, and pairwise
coprimality.  These counterexamples retire the exact one-move policies only.
Gate VF permits arbitrary additional five-term moves and remains open.  The
abc conjecture also remains open.

The Lean companion
\texttt{QuadraticVeronesePeeling20260902.lean} formalizes the exact domain,
integer transport, valuation-layer identities, and all three quantified
fixed-policy no-go theorems.  Its separate axiom audit covers all 53 public
theorem declarations.  A deterministic computation checks 20,000
consecutive parameters and a power-of-two subsequence through $2^{18}$.
